\paragraph{Calibration is not informativeness.} A predictive distribution can be calibrated, sharp enough to rival a naive baseline and correlated with its own errors without using its input. In EEG-to-fMRI synthesis this is the default: the target is dominated by anatomy that is shared across subjects, and anatomical variability determines both where errors occur and where any reasonable model is uncertain. All eleven methods learn this structure, and so does the EEG-free template, which none of them beats significantly. Their uncertainty is a fixed anatomical map that does not move with the EEG. Had we reported only NLL, PICP, Spearman and AUSE, as is common in UQ papers, the ensemble, MC dropout and the stochastic spectral decoder would have looked like successful uncertainty quantification.

\paragraph{Heteroscedastic training can blind a synthesizer to its input.} The positive control shows that the problem is not only the data. A heteroscedastic likelihood lets a network explain a residual either by moving its mean or by widening its variance. When the input is weakly informative and the learning rate is high, widening wins, and the mean stops fitting \citep{seitzer2022pitfalls}. In our setting this turned a backbone that can recover an injected EEG signal ($r=0.48$ with a point loss) into one that cannot ($r\approx0$), while its NLL, coverage and sparsification remained respectable. Every UQ method built on that solution (heteroscedastic, Laplace, SWAG, ensembles) inherits the blindness. This is the strongest argument for positive controls in UQ evaluation: a method's calibration says nothing about whether it can see a signal that is there.

\paragraph{What the EEG predicts, and what it does not.} Models that can see an EEG-coupled signal still find little in the real data. At the voxel level, the point-loss decoder and a linear probe on the same input exceed the time-shift null only through the shared task timing; neither exceeds the subject-swap null. At the region level, a lagged linear model predicts motor BOLD across subjects as convincingly as any published result. But heart rate predicts the same region as well, the EEG adds nothing to heart rate and motion across subjects and at most a small within-subject increment that does not survive the ECG, and the EEG still contains cardiac signal after pulse-artefact correction. The simplest reading is that the region-level prediction is physiological: heart-rate fluctuations drive BOLD changes through cerebral blood flow, and they enter the EEG through residual pulse artefacts and autonomic covariation. This does not mean that EEG--BOLD coupling is absent. Alpha and mu/beta power covary with BOLD in specific networks \citep{goldman2002simultaneous,laufs2003eeg}. It means that in this dataset a predictive model cannot separate that coupling from physiology, and that a surrogate test over the BOLD, including the careful test of Mehta et al., cannot either. Region-level synthesis models report cross-subject correlations of a similar size on other datasets \citep{li2024neurobolt}; we suggest that they report heart-rate and motion models, and nested comparisons, alongside them.

\paragraph{Practices.} The protocol is general, and we suggest the following for EEG-to-fMRI synthesis and similar cross-modal tasks:
\begin{enumerate}
	\item report proper scores (NLL, CRPS) of an EEG-free predictive distribution next to every model, including a stimulus-locked one when the task timing is shared;
	\item evaluate every trained model under a temporal-shift null and a subject-swap null, and claim EEG-specific information only above the swap null;
	\item test whether the predicted uncertainty changes with the input, not only whether it is calibrated;
	\item include a positive control: inject a known input-coupled signal and show that the model, trained exactly as reported, recovers it;
	\item fit physiological confound models (heart rate from the ECG, head motion) with the same pipeline, and test the EEG's increment over them in nested models;
	\item split by subject, calibrate conformal intervals over many subjects (cross-conformally if needed), and report coverage per subject;
	\item check the EEG--fMRI alignment from the scanner markers and the spectral bins against the Nyquist frequency, register the volumes from the scanner's world coordinates when the field of view is partial, and check every registration visually;
	\item report effect sizes with confidence intervals and a smallest effect of interest for every claim of absence.
\end{enumerate}

\paragraph{Limitations.} The analyses use two motor-imagery datasets from one laboratory, with 17 and 9 subjects and block designs; the region-level confound result is weaker on the smaller, noisier-ECG cohort. Subject-level tests with $n=17$ have limited power, but the equivalence tests bound the effects we could have missed (Table~\ref{tab:effects}). The positive control injects a linear, spatially fixed signal; real EEG--BOLD coupling may be nonlinear and spatially specific, and we test sensitivity, not the full range of plausible couplings. The backbone's architecture was fixed by the original neural architecture search, and the UQ methods share its searched learning rate. A lower learning rate rescues the heteroscedastic head on the positive control, but we did not retune and rerun every method, and the conclusions about each method apply to it as configured. Heart rate was derived from the in-scanner ECG, which contains gradient residuals; at most 9\% of RR intervals per subject were implausible and interpolated. Respiration was not recorded and could account for part of what we attribute to heart rate. Both arms normalize each subject with its own run (transductive normalization), as is standard in this literature; the EEG-free references are scored on the same normalized data. Finally, the registered-volume analyses used a corrected registration, and the native-space analyses inherit the geometry of the original backbones (Section~\ref{sec:res:registered}).
